\documentclass[11pt]{article}
\usepackage[a4paper,margin=1in]{geometry}
\usepackage{amsmath,amssymb,mathtools,bm}
\usepackage{enumitem,booktabs,array}
\usepackage{tcolorbox}
\usepackage{hyperref}
\usepackage[protrusion=true,expansion=false]{microtype}
\hypersetup{colorlinks=true,linkcolor=blue,urlcolor=blue,citecolor=blue}
\setlist[itemize]{topsep=2pt,itemsep=2pt}
\setlist[enumerate]{topsep=2pt,itemsep=2pt}
\newtcolorbox{keybox}{colback=blue!3!white,colframe=blue!40!black,boxrule=0.4pt,arc=1mm}
\newtcolorbox{alertbox}{colback=red!3!white,colframe=red!40!black,boxrule=0.4pt,arc=1mm}
\newtcolorbox{answerbox}{colback=green!3!white,colframe=green!40!black,boxrule=0.4pt,arc=1mm}
\newtcolorbox{drillbox}{colback=yellow!5!white,colframe=orange!60!black,boxrule=0.4pt,arc=1mm}
\newcommand{\EV}{\mathbb{E}}
\newcommand{\Var}{\mathrm{Var}}
\newcommand{\Cov}{\mathrm{Cov}}
\newcommand{\blank}[1]{\underline{\hspace{#1}}}

\title{E10-04: Ewens, Nanda \& Rhodes-Kropf (2018, JFE) \\
\large ``Cost of Experimentation and the Evolution of Venture Capital'' --- Active Recall Workbook}
\author{Empirical Corporate Finance II --- Exam Prep}
\date{\today}
\begin{document}\maketitle

\begin{keybox}
\textbf{Paper in one sentence.} ENRK argue that a technology-induced fall in the cost of experimentation (cloud computing, open-source tools, lean-startup methods) transformed VC financing: more startups attempt bold early-stage experiments, financing rounds are smaller and more staged, success rates fall but expected returns rise, and VC portfolios shift toward many small bets with richer-options structure.

\textbf{Placement.} Important paper explaining the modern VC-industry structure (seed vs.\ late-stage). Connects innovation economics to financing contracts and real-options theory.
\end{keybox}

\section*{Round 1: Complete Walk-Through}

\subsection*{1.1 Research question}
In the 2000s, cloud computing and open-source tooling sharply reduced the cost of launching a tech startup. What effect did this have on the structure of VC investment and on startup outcomes?

\subsection*{1.2 Theory}
Lower experimentation cost $\Rightarrow$ cheaper to gain information. Real-options view predicts:
\begin{itemize}
\item More startups attempt radical experiments.
\item Investors fund smaller seed rounds to learn before committing.
\item More startups fail at an early stage (``fail fast'').
\item Those that succeed scale more aggressively (larger follow-on rounds).
\end{itemize}

\subsection*{1.3 Data}
\begin{itemize}
\item VentureSource + AngelList + Crunchbase US seed-stage IT companies, 1995--2014.
\item Compared across pre-cloud (1995--2004) and post-cloud (2005--2014) periods.
\item Financing-round size, startup outcomes (IPO, acquisition, failure), follow-on funding.
\end{itemize}

\subsection*{1.4 Headline results}
\begin{itemize}
\item Seed investment counts rise $\sim$2--3x; seed round sizes fall.
\item Failure rates rise at seed stage; conditional-on-survival revenue growth higher.
\item Portfolio returns are more skewed: many more total failures, but successes larger.
\item VC industry bifurcates: seed VCs, growth VCs.
\item Consistent with the real-options, lower-experimentation-cost hypothesis.
\end{itemize}

\subsection*{1.5 Mechanism}
\begin{itemize}
\item Technology reduces the marginal cost of testing startup concepts.
\item Investors can gather information cheaply through staged experiments.
\item Portfolio optimization favors a larger, more diversified book of small seed bets.
\end{itemize}

\subsection*{1.6 Implications}
\begin{itemize}
\item VC industry structure co-evolves with technology.
\item Policymakers should view startup failure rates in the context of cheaper experimentation.
\item Innovation economics must incorporate experimentation costs explicitly.
\end{itemize}

\subsection*{1.7 Caveats}
\begin{alertbox}
\begin{itemize}
\item Cross-period identification, not a clean natural experiment.
\item Measurement of ``experimentation cost'' is indirect.
\item Industry-specific shocks could confound.
\end{itemize}
\end{alertbox}

\section*{Round 2: Level 1 Fill-in}
\begin{enumerate}
\item Technology shock: \blank{2cm} computing + open-source tools.
\item Pre-cloud: \blank{1cm}--\blank{1cm}; post-cloud: \blank{1cm}--\blank{1cm}.
\item Seed deal count: $\times$\blank{1cm}--\blank{1cm}.
\item Round size trend: \blank{2cm}.
\item Portfolio distribution: more \blank{2cm}.
\end{enumerate}

\section*{Round 3: Level 2 Prompts}
\begin{enumerate}
\item State the theory linking experimentation cost to VC structure.
\item Describe the data and period comparisons.
\item Report main findings on seed rounds, failure, and success.
\item Discuss mechanism: real-options optimization.
\item Policy and industry implications.
\end{enumerate}

\section*{Round 4: Minimal Scaffolding}
From a blank page: (i) theory, (ii) data, (iii) findings, (iv) mechanism, (v) implications.

\section*{Round 5: Blank Page}
Essay: cost of experimentation and the modern VC industry.

\vspace{6pt}
\begin{drillbox}
\textbf{Speed Drill.} (1) Shock. (2) Seed-count change. (3) Round-size trend. (4) Failure pattern. (5) Industry implication.
\end{drillbox}

\section*{Quick Reference \& Answer Key}
\begin{answerbox}
\begin{itemize}
\item \textbf{Shock.} Cloud + open-source reduce exp cost.
\item \textbf{Seed.} More, smaller rounds.
\item \textbf{Failure.} Rises early, successes bigger.
\item \textbf{Structure.} Seed vs.\ growth VC bifurcation.
\end{itemize}
\end{answerbox}

\begin{keybox}
\textbf{30-second exam pitch.} Ewens, Nanda \& Rhodes-Kropf (2018) argue that the 2000s decline in startup-experimentation costs (cloud computing, open-source, lean-startup methods) transformed VC financing. Using VentureSource/AngelList/Crunchbase data 1995--2014, they compare pre- and post-cloud periods: seed-stage deal counts rise 2--3x, round sizes fall, failure rates at seed rise, but conditional-on-survival growth accelerates; portfolio returns skew more sharply, with many more failures offset by larger successes. The VC industry bifurcates into seed and growth funds. This matches a real-options model where lower experimentation cost favors many small, staged bets with richer learning value. The paper provides a technology-driven explanation for the structure of modern VC, with implications for innovation economics and startup policy: failure-rate increases do not imply declining efficiency --- they signal cheaper information acquisition.
\end{keybox}

\end{document}
